\documentclass[aps,prd,twocolumn,superscriptaddress,nofootinbib,longbibliography]{revtex4-2}

\usepackage[utf8]{inputenc}
\usepackage[T1]{fontenc}
\usepackage{amsmath,amssymb,bm}
\usepackage{graphicx}
\usepackage{booktabs}
\usepackage{array}
\usepackage[colorlinks=true,linkcolor=blue,citecolor=blue,urlcolor=blue]{hyperref}

\newcommand{\dd}{\mathrm{d}}
\newcommand{\Lag}{\mathcal{L}}
\newcommand{\calS}{\mathcal{S}}
\newcommand{\calC}{\mathcal{C}}
\newcommand{\calD}{\mathcal{D}}
\newcommand{\Mp}{M_{\rm p}}
\newcommand{\Lam}{\Lambda}
\newcommand{\LCDM}{\Lambda{\rm CDM}}
\newcommand{\DES}{\mathrm{DES}}
\newcommand{\DESI}{\mathrm{DESI}}
\newcommand{\SN}{\mathrm{SN}}
\newcommand{\BAO}{\mathrm{BAO}}
\newcommand{\varphione}{\varphi_1}
\newcommand{\varphitwo}{\varphi_2}
\newcommand{\chiT}{\chi^2_{\rm tot}}

\begin{document}

\preprint{AI-MG-v5.2.2-agent}

\title{Autonomous Search for Viable Two-Scalar Modified Gravity\\
from Cosmological Background Data}

\author{Cooper Niu}
\affiliation{Department of Physics, Brown University, Providence, Rhode Island 02912, USA}

\author{Wenrong Sun}
\affiliation{Department of Physics, Brown University, Providence, Rhode Island 02912, USA}

\date{\today}

\begin{abstract}
We present a physics-first framework for the autonomous discovery of two-scalar modified-gravity theories from late-time cosmological data.  Candidate theories are generated as covariant, dimension-balanced Lagrangians built from curvature tensors, covariant derivatives, and two scalar fields.  They are then subjected to deterministic physical validation: the FLRW equations are derived symbolically, higher-derivative terms are removed by ghost-free coefficient relations, the lapse equation is retained as a modified Friedmann constraint, and the surviving systems are integrated and compared with DES type-Ia supernova distances and DESI Data Release 2 baryon-acoustic-oscillation measurements.  The artificial-intelligence component is used only to navigate the large discrete theory space and propose parameter samples; theoretical consistency and empirical scoring are decided by the physics pipeline.  We emphasize the role of temporal diffeomorphism invariance in the minisuperspace theory: even when a modified-gravity Lagrangian contains derivatives of the lapse before simplification, the residual time-reparametrization Noether identity implies that, in the ghost-free FLRW sector, the gauge-fixed lapse equation is a first-order constraint rather than an independent dynamical equation, and this constraint is preserved by the evolution equations.  Recent v5.2.2-agent scans find background solutions with $\chiT\simeq 1638$ for DES-SN plus DESI-DR2 BAO.  The best families are typically Einstein--Hilbert anchored scalar sectors, while a subset contains weak, moderate, or active curvature corrections.  These results identify promising background branches, not final fully viable gravitational theories; perturbation stability, structure growth, gravitational-wave constraints, and local tests remain necessary future validation layers.
\end{abstract}

\maketitle

\section{Introduction}
\label{sec:introduction}

The observed acceleration of the Universe is one of the clearest indications that either the cosmic energy budget or the gravitational sector is not yet fully understood.  The minimal $\LCDM$ model describes current background data with impressive efficiency, but it leaves the cosmological constant problem and the microscopic origin of dark energy unresolved.  Modified gravity and scalar-tensor theories provide a systematic way to test whether late-time acceleration can arise from additional gravitational degrees of freedom rather than from a rigid vacuum energy \cite{Clifton2012,SotiriouFaraoni2010,Horndeski1974}.

The difficulty is not the absence of possible theories.  It is the opposite: the space of covariant Lagrangians built from curvature invariants, derivatives, and scalar fields is combinatorially large.  Most expressions are physically unacceptable.  Generic higher-derivative theories contain Ostrogradsky instabilities \cite{Ostrogradsky1850}; many theories fail to admit a regular cosmological initial-value problem; and many stable background branches do not fit the data.  Historically, theorists have navigated this landscape with symmetry, naturalness, simplicity, and experience.  These principles remain essential, but they are incomplete search strategies in a high-dimensional theory space.

Recent DES supernova and DESI baryon-acoustic-oscillation measurements provide a sharp background-expansion target for such searches \cite{DES2024SN,DESI2025BAO}.  At the same time, recent work on autonomous theory discovery has shown that physics can be encoded as a formal language, so that artificial intelligence proposes only well-formed theories while deterministic computations impose first-principle constraints and evaluate data likelihoods \cite{Alexander2026Albert}.  We adopt the same philosophy for modified gravity.  The agent is not asked to hallucinate equations or to judge its own success.  It proposes covariant operator combinations and numerical samples inside a fixed grammar.  The theory is then accepted or rejected by an explicit symbolic and numerical pipeline.

This paper describes the physical content of the v5.2.2-agent modified-gravity framework.  The emphasis is deliberately different from a software manual.  The implementation details are documented in the public code; here we focus on the theory space, the ghost-free FLRW reduction, the Hamiltonian-type constraint, the construction of a dimensionless cosmological initial-value problem, and the interpretation of the resulting DES-SN and DESI-DR2 BAO fits.

\section{Covariant Theory Space}
\label{sec:theory-space}

\subsection{Dimension-balanced effective Lagrangians}
\label{subsec:dimension-balanced}

We consider a broad EFT-like space of metric theories containing two scalar fields, $\phi_1$ and $\phi_2$.  A generated theory is written as
\begin{equation}
\calS=\int \dd^4x\,\sqrt{-g}\,\Lag ,
\label{eq:action}
\end{equation}
with
\begin{equation}
\Lag=
\Mp^2\sum_{I\in G\cup N}
c_I\Lam^{2-d_I}O_I
+
\sum_{J\in S}c_J\Lam^{4-d_J}S_J
+
\Lag_{\rm fixed}.
\label{eq:lagrangian}
\end{equation}
Here $G$ denotes pure-curvature invariants, $N$ non-minimal scalar-curvature invariants, and $S$ pure-scalar generated invariants.  The integer $d_I$ is the mass dimension of the corresponding covariant scalar.  The factors of $\Mp^2$ and $\Lam$ ensure that each term in $\sqrt{-g}\Lag$ has the correct mass dimension.  Thus a candidate is not specified by an operator label alone; its sector, coefficient, cutoff power, and Planck-mass prefactor are part of the physical theory.

The candidate building blocks include
\begin{equation}
R,\qquad R_{\mu\nu},\qquad R^\rho{}_{\sigma\mu\nu},\qquad
\nabla_\mu,\qquad \phi_1,\qquad \phi_2 .
\end{equation}
Examples of generated invariants are
\begin{equation}
R^2,\qquad
R_{\mu\nu}R^{\mu\nu},\qquad
\phi_1^2R,\qquad
\nabla_\mu\phi_1\nabla_\nu\phi_1R^{\mu\nu},\qquad
\Box\phi_1\,R .
\end{equation}
The scalar configurations are
\begin{equation}
\hbox{none},\qquad \phi_1\hbox{ only},\qquad \phi_2\hbox{ only},\qquad \phi_1+\phi_2 .
\end{equation}
When a scalar is active, the fixed matter sector contains a canonical kinetic term and either a quadratic or an axion-like potential:
\begin{align}
\Lag_{\rm fixed}^{(i)}
&=-\frac12 g^{\mu\nu}\nabla_\mu\phi_i\nabla_\nu\phi_i - V_i(\phi_i),\\
V_i(\phi_i)
&=\frac12 m_i^2\phi_i^2
\quad{\rm or}\quad
m_i^2f_i^2\left[1-\cos\left(\frac{\phi_i}{f_i}\right)\right].
\label{eq:fixed-sector}
\end{align}

The plain $R$ term, when present, is treated as an Einstein--Hilbert anchor with coefficient fixed to $c_R=1/2$ in the numerical comparison:
\begin{equation}
\Lag_{\rm EH}=\frac{\Mp^2}{2}R.
\label{eq:EH}
\end{equation}
This common normalization allows meaningful comparison of higher-curvature and scalar-curvature corrections across candidate theories.

\subsection{Why two scalar normalizations?}
\label{subsec:two-scalars}

The two scalar fields play different roles in the current theory grammar.  The first scalar is naturally measured in Planck units,
\begin{equation}
\phi_1=\Mp\varphione,
\end{equation}
while the second is measured in cutoff units,
\begin{equation}
\phi_2=\Lam\varphitwo .
\end{equation}
This separation allows the search to include both Planck-normalized scalar-tensor sectors and cutoff-normalized EFT scalar structures.  The distinction matters because a term such as $\phi_2 R$ carries a different physical size from $\phi_1 R$, even when their dimensionless field amplitudes are comparable.  The Solver therefore tracks the cutoff ratio
\begin{equation}
\lambda_H=\frac{\Lam}{H_0}
\end{equation}
and enforces an EFT consistency condition along accepted trajectories,
\begin{equation}
\max_x\frac{|H(x)|}{\Lam}=\max_x\frac{|h(x)|}{\lambda_H}<0.1.
\label{eq:eft-cut}
\end{equation}

\subsection{Free parameters and scan ranges}
\label{subsec:parameter-priors}

The autonomous search is meaningful only after specifying which continuous quantities are allowed to vary.  Table~\ref{tab:parameter-priors} summarizes the default cold-start parameter ranges used in the v5.2.2-agent scans.  These are proposal priors, not physical acceptance conditions.  A sample drawn from these ranges is still rejected if it fails the Friedmann constraint, violates the EFT cut, produces an unstable trajectory, or gives no acceptable real initial-velocity root.

\begin{table*}[t]
\centering
\begin{ruledtabular}
\begin{tabular}{lllll}
Parameter & Dimensionless definition & Applies to & Default range & Sampling scale \\
\hline
$H_0$ & $H_0\,[{\rm km\,s^{-1}\,Mpc^{-1}}]$ & all sectors & $58$--$74$ & linear \\
$\Omega_{m0}$ & $\rho_{m0}/(3\Mp^2H_0^2)$ & all sectors & $0.22$--$0.40$ & linear \\
$\lambda_H$ & $\Lam/H_0$ & all sectors & $5.0\times10^2$--$8.0\times10^4$ & log \\
$\varphi_{1,0}$ & $\phi_1(t_0)/\Mp$ & $\phi_1$ sectors & $0.15$--$4.5$ & linear \\
$\mu_1$ & $m_1/H_0$ & $\phi_1$ sectors & $0.04$--$4.0$ & log \\
$\bar f_1$ & $f_1/\Mp$ & $\phi_1$ axion sectors & $0.05$--$30$ & log \\
$\varphi_{2,0}$ & $\phi_2(t_0)/\Lam$ & $\phi_2$ sectors & $-3.0$--$3.0$ & linear \\
$q_{2,0}$ & $\dd\varphi_2/\dd\ln a|_{0}$ & two-scalar sectors & $-5.0$--$5.0$ & linear \\
$\mu_2$ & $m_2/\Lam$ & $\phi_2$ sectors & $10^{-4}$--$1$ & log \\
$\bar f_2$ & $f_2/\Lam$ & $\phi_2$ axion sectors & $0.05$--$50$ & log \\
$c_I$ & non-EH operator coefficients & generated operators & $10^{-10}$--$10^{-4}$ in $|c_I|$ & log, both signs \\
\end{tabular}
\end{ruledtabular}
\caption{Default continuous proposal ranges.  The Einstein--Hilbert coefficient is not scanned: when the plain $R$ anchor is present it is fixed to $c_R=1/2$.  In one-scalar models the initial scalar velocity is solved from the modified Friedmann constraint rather than independently sampled.  In two-scalar models $q_{2,0}$ is sampled and the constraint solves for the remaining scalar velocity.}
\label{tab:parameter-priors}
\end{table*}

For higher-dimension curvature corrections, a naive coefficient prior can make the correction effectively invisible because of powers of $\Lam$.  The scans therefore also use an effective-size diagnostic and, in some proposal modes, choose coefficient magnitudes to target weak, moderate, or active fractional corrections relative to the Einstein--Hilbert term.  The diagnostic range is
\begin{equation}
\Delta_i\simeq
\left|\frac{c_i}{c_{\rm EH}}\right|\lambda_H^{p_i}{\cal F}_i ,
\label{eq:delta-prior}
\end{equation}
with weak, moderate, and active targets corresponding approximately to $10^{-6}$--$10^{-4}$, $10^{-4}$--$10^{-2}$, and $10^{-2}$--$5\times10^{-2}$, respectively.  This prescription is a search prior only; the final ranking remains the real data value of $\chiT$.

\section{FLRW Reduction and Ghost-Free Filtering}
\label{sec:ghost}

\subsection{Background equations}
\label{subsec:background-eom}

The first physical filter is applied after reducing the covariant theory on the spatially flat FLRW metric
\begin{equation}
\dd s^2=-N(t)^2\dd t^2+a(t)^2\delta_{ij}\dd x^i\dd x^j .
\label{eq:flrw}
\end{equation}
The minisuperspace variation gives
\begin{equation}
\delta S=\int \dd t\left(E_N\delta N+E_a\delta a+\sum_iE_{\phi_i}\delta\phi_i\right),
\label{eq:variation}
\end{equation}
where the inactive scalar equations are omitted.  After variation, the equations are evaluated in cosmic-time gauge, $N=1$:
\begin{equation}
E_N=0,\qquad E_a=0,\qquad E_{\phi_i}=0 .
\label{eq:eom-set}
\end{equation}
The lapse equation is not discarded.  It is the modified Friedmann equation, while $E_a$ and $E_{\phi_i}$ determine the dynamical evolution.

\subsection{Ostrogradsky danger and background ghost-free relations}
\label{subsec:ostro}

Higher-curvature and derivative scalar-curvature operators generically produce equations containing third or higher time derivatives of $a(t)$ and $\phi_i(t)$.  If such derivatives represent independent initial data, the theory contains an Ostrogradsky mode: the Hamiltonian is unbounded and the background branch is physically unacceptable \cite{Ostrogradsky1850}.  The pipeline therefore treats the derivative order of the FLRW equations as a hard physical constraint.

Operationally, each candidate Lagrangian is inspected for higher-derivative structures.  The EOM stage solves for coefficient relations that cancel all forbidden higher-derivative terms in the gauge-fixed background equations.  A theory may survive because a special combination of operators is degenerate on FLRW, or it may fail because no such relation exists.  Passing this filter does not prove full nonlinear perturbative stability, but it removes the most immediate background Ostrogradsky instability and defines the theory sector that is eligible for numerical cosmology.

This is analogous in spirit to the way Horndeski and degenerate scalar-tensor theories avoid extra propagating ghost modes through special derivative structures \cite{Horndeski1974,Langlois2016DHOST}.  Our search is less restrictive at the proposal stage: it allows a broad grammar, then lets the symbolic EOM filter identify which combinations have acceptable FLRW derivative order.

\section{The Modified Friedmann Constraint}
\label{sec:hamiltonian}

\subsection{Residual time diffeomorphism identity}
\label{subsec:noether}

In ordinary GR, the FLRW lapse $N$ appears without time derivatives after the Einstein--Hilbert boundary term is handled.  The momentum conjugate to $N$ vanishes, and preserving this primary constraint gives the Hamiltonian constraint, which reduces to the Friedmann equation.  In a raw modified-gravity minisuperspace Lagrangian, however, $N$ may appear with derivatives before integrations by parts or before imposing ghost-free relations.  Therefore the simple GR statement ``the lapse is a Lagrange multiplier'' is not always the cleanest way to identify the constraint.

The more robust argument comes from residual time-reparametrization invariance.  Under
\begin{equation}
t\rightarrow t-\xi(t),
\end{equation}
the minisuperspace variables transform as
\begin{equation}
\delta_\xi a=\xi\dot a,\qquad
\delta_\xi\phi_i=\xi\dot\phi_i,\qquad
\delta_\xi N=(N\xi)^{\cdot}.
\label{eq:time-diffeo}
\end{equation}
Invariance of the action under this transformation gives the off-shell Noether identity
\begin{equation}
N\dot E_N-\dot a\,E_a-\sum_i\dot\phi_i\,E_{\phi_i}=0.
\label{eq:noether}
\end{equation}
This is the minisuperspace form of the Bianchi identity.  It holds before imposing the equations of motion.

\subsection{Why the lapse equation is a constraint}
\label{subsec:constraint-proof}

After gauge fixing to $N=1$, Eq.~\eqref{eq:noether} becomes
\begin{equation}
\dot E_N=\dot a\,E_a+\sum_i\dot\phi_i\,E_{\phi_i}.
\label{eq:noether-gauge}
\end{equation}
Suppose the ghost-free relations have made the dynamical equations $E_a=0$ and $E_{\phi_i}=0$ at most second order in time derivatives.  Then the right-hand side of Eq.~\eqref{eq:noether-gauge} contains no independent third derivatives.

If the gauge-fixed lapse equation $E_N|_{N=1}$ contained an acceleration, schematically
\begin{equation}
E_N|_{N=1}=C(a,\dot a,\ddot a,\phi_i,\dot\phi_i,\ddot\phi_i),
\end{equation}
then its time derivative would contain independent third derivatives:
\begin{equation}
\dot E_N|_{N=1}
\supset
\frac{\partial C}{\partial \ddot a}a^{(3)}
+\sum_i
\frac{\partial C}{\partial \ddot\phi_i}\phi_i^{(3)}.
\end{equation}
These terms cannot be matched by the right-hand side of Eq.~\eqref{eq:noether-gauge}.  Since the identity is off shell and holds for arbitrary field profiles, the coefficients of all independent third derivatives must vanish.  Therefore
\begin{equation}
E_N|_{N=1}=C_N(a,\dot a,\phi_i,\dot\phi_i),
\label{eq:first-order-constraint}
\end{equation}
and the lapse equation is a first-order constraint on FLRW initial data, not an independent second-order dynamical equation.

\subsection{Constraint preservation}
\label{subsec:constraint-preservation}

The same identity guarantees propagation of the constraint.  On solutions of the dynamical equations,
\begin{equation}
E_a=0,\qquad E_{\phi_i}=0,
\end{equation}
Eq.~\eqref{eq:noether-gauge} implies
\begin{equation}
\dot E_N=0.
\end{equation}
Thus if the initial data satisfy
\begin{equation}
E_N(t_0)=0,
\end{equation}
then the background evolution preserves
\begin{equation}
E_N(t)=0.
\end{equation}
This justifies using the lapse equation to solve for one initial velocity and then monitoring the residual constraint throughout numerical integration.

This argument should not be overstated.  It is a background minisuperspace result, not a full ADM/Dirac proof of the complete four-dimensional constraint algebra \cite{Dirac1958,ADM1962}.  What it establishes is precisely what the cosmological Solver needs: in the ghost-free FLRW sector, the lapse equation is the modified Friedmann constraint and is preserved by the accepted dynamical equations.

\section{Dimensionless Cosmological Initial-Value Problem}
\label{sec:ivp}

The numerical evolution uses e-fold time
\begin{equation}
x=\ln a=-\ln(1+z),
\qquad
h=\frac{H}{H_0},
\label{eq:xh}
\end{equation}
so that
\begin{equation}
\frac{\dd}{\dd t}=H_0h\frac{\dd}{\dd x}.
\label{eq:chain-rule}
\end{equation}
The relevant substitutions include
\begin{align}
\dot a&=H_0e^xh,
&
\ddot a&=H_0^2e^x(h^2+hh_x),\\
\dot\phi_1&=\Mp H_0h q_1,
&
\ddot\phi_1&=\Mp H_0^2h(h_xq_1+hq_{1x}),\\
\dot\phi_2&=H_0^2\lambda_H h q_2,
&
\ddot\phi_2&=H_0^3\lambda_H h(h_xq_2+hq_{2x}),
\label{eq:rescalings}
\end{align}
where
\begin{equation}
q_i=\frac{\dd\varphi_i}{\dd x}.
\end{equation}
Masses and decay constants are written as
\begin{equation}
m_1=H_0\mu_1,\qquad
m_2=\Lam\mu_2,\qquad
f_1=\Mp \bar f_1,\qquad
f_2=\Lam \bar f_2 .
\end{equation}

After these substitutions, the Builder forms dimensionless equations
\begin{equation}
C_N=0,\qquad C_a=0,\qquad C_{\phi_i}=0.
\end{equation}
The constraint $C_N=0$ is imposed at $x=0$ with $h(0)=1$.  The dynamical equations are solved algebraically for the highest e-fold derivatives,
\begin{equation}
\bm u=(h_x,q_{1x},q_{2x})^T.
\end{equation}
When a regular branch exists, the result is an explicit first-order system,
\begin{equation}
\frac{\dd}{\dd x}
\begin{pmatrix}
h\\ \varphione\\ q_1\\ \varphitwo\\ q_2
\end{pmatrix}
=
\begin{pmatrix}
F_h\\ q_1\\ F_1\\ q_2\\ F_2
\end{pmatrix},
\label{eq:first-order}
\end{equation}
with absent scalar variables omitted.  The trajectory is rejected if the constraint drifts beyond the production tolerance,
\begin{equation}
\max_x|C_N(x)|>10^{-4},
\label{eq:constraint-cut}
\end{equation}
or if the solution violates the EFT condition Eq.~\eqref{eq:eft-cut}.

\section{Cosmological Observables}
\label{sec:observables}

Each accepted solution gives a Hubble history $H(z)=H_0h(z)$.  The supernova distance modulus is computed from
\begin{equation}
D_L(z)=(1+z_{\rm hel})\frac{c}{H_0}
\int_0^{z_{\rm HD}}\frac{\dd z'}{h(z')},
\label{eq:luminosity-distance}
\end{equation}
with an analytically profiled absolute-magnitude offset.  The supernova likelihood therefore mainly constrains the relative shape of the distance-redshift relation.

The BAO observables are computed from the same background solution:
\begin{align}
\calD_M(z)&=\int_0^z\frac{\dd z'}{h(z')},\\
\calD_H(z)&=\frac{1}{h(z)},\\
\calD_V(z)&=\left[\frac{z\calD_M(z)^2}{h(z)}\right]^{1/3}.
\label{eq:bao}
\end{align}
In the current background-only analysis, the BAO scale is profiled as a nuisance parameter.  Thus the DESI comparison should be interpreted as a relative-distance test of the late-time expansion history, not as a complete early-universe calculation of the sound horizon in each modified-gravity theory.

The primary score is
\begin{equation}
\chiT=\chi^2_{\DES\text{-}\SN}+\chi^2_{\DESI\text{-DR2}\,\BAO}.
\label{eq:chi2}
\end{equation}
Models are ranked by $\chiT$, with information criteria recorded only as secondary diagnostics \cite{Akaike1974,Schwarz1978}.

\section{Autonomous Search Strategy}
\label{sec:search}

The autonomous component follows the same principle as formal-language approaches to particle-physics discovery: generate theories inside a grammar, enforce hard theoretical constraints, compute observables with deterministic code, and reward agreement with data \cite{Alexander2026Albert}.  In the present work the grammar is a modified-gravity Lagrangian grammar rather than a particle-physics gauge grammar.  The hard constraints are covariance of the generated operators, dimension balancing, background ghost-free relations, existence of a regular IVP, Friedmann-constraint preservation, and EFT validity.  The empirical reward is the Solver-written $\chiT$.

The important point is that the search policy is not part of the physics claim.  Different proposal policies can explore the same theory space.  A candidate becomes physically interesting only when it passes the deterministic validation chain and produces an accepted cosmological trajectory.  For that reason we report accepted theories by their complete Lagrangians, coefficient values, scalar initial conditions, constraint residuals, and data-fit metrics rather than by agent prompts or internal training details.

\section{Representative Results}
\label{sec:results}

Recent v5.2.2-agent scans provide a useful snapshot of the current search.  They are not combined into a posterior distribution, since the proposal policies and run budgets differ.  They should instead be read as validated examples of what the pipeline can find and as diagnostics for where the theory space is difficult.

\begin{table*}[t]
\centering
\begin{ruledtabular}
\begin{tabular}{lrrrrl}
Run & Attempts & Accepted & Families & Best $\chiT$ & Main physical lesson \\
\hline
\texttt{graph\_agent\_batch\_1600\_llm\_v3} & 880 & 309 & 49 & 1637.966 & Competitive EH-anchored scalar sectors and some active corrections \\
\texttt{graph\_agent\_batch\_208\_split\_v4} & 208 & 70 & 20 & 1638.043 & EH-anchored track strongly outperforms non-EH-only track \\
\texttt{graph\_agent\_oscar\_llm\_strict\_500\_nomem} & 561 & 310 & 18 & 1638.835 & Strict no-memory run rediscovers simple $\phi_1$ sectors \\
\end{tabular}
\end{ruledtabular}
\caption{Representative recent validated agent runs.  ``Families'' denotes the number of best-per-family representatives shown in the human-facing reports; the machine-readable reports keep all accepted attempts.}
\label{tab:results}
\end{table*}

The best representative in the 1600-attempt report has
\begin{equation}
\chiT=1637.966,
\qquad
H_0=65.599~{\rm km\,s^{-1}\,Mpc^{-1}},
\qquad
\Omega_{m0}=0.2727,
\end{equation}
with $\lambda_H=3916.603$.  Its operator content is an Einstein--Hilbert anchor plus a derivative curvature correction,
\begin{equation}
R,\qquad \nabla^cR^{ab}\nabla_cR_{ab},
\end{equation}
in a two-scalar axion/axion sector.  The effective correction diagnostic classifies the curvature correction as moderate, with relative size $\Delta_{\rm max}=4.47\times10^{-4}$.

The same scan also finds Einstein--Hilbert plus scalar baselines with nearly identical background quality.  For example, a $\phi_1$ axion model with only the $R$ anchor reaches
\begin{equation}
\chiT=1638.021.
\end{equation}
This is physically important: the data reward smooth scalar dark-energy expansion histories, and a low background $\chi^2$ by itself does not prove that an additional curvature correction is dynamically important.

The split-track run makes this point even more clearly.  Its best EH-anchored representative is a $\phi_1$ quadratic model with only the $R$ operator and
\begin{equation}
\chiT=1638.043.
\end{equation}
Several nearby models include corrections such as $\phi_1\Box R$, $R^2$, $\nabla\phi_1\nabla\phi_1R_{\mu\nu}$, and $\phi_1R_{\mu\nu\rho\sigma}R^{\mu\nu\rho\sigma}$, with weak-to-active effective sizes.  In the same run, the best non-EH-only and $\phi_2$-dominated representatives have $\chiT\simeq9434$, indicating that these branches are much less successful under the current priors and likelihood.  This should be interpreted as a current-loop empirical signal, not as a theorem excluding those sectors.

\subsection{Effective modified-gravity signal}
\label{subsec:effective-signal}

To distinguish scalar dark-energy baselines from genuinely active modified-gravity corrections, accepted models are assigned an effective correction diagnostic.  For each non-EH curvature correction, we estimate
\begin{equation}
\Delta_i\simeq
\left|\frac{c_i}{c_{\rm EH}}\right|
\lambda_H^{p_i}{\cal F}_i,
\label{eq:delta}
\end{equation}
where $p_i$ is the validated cutoff power and ${\cal F}_i$ is a dimensionless field factor inferred from the accepted sample.  Pure scalar self-interactions are not counted as modified-gravity corrections in this diagnostic.  The maximum $\Delta_{\rm max}$ is classified as active ($\ge10^{-2}$), moderate ($10^{-4}$--$10^{-2}$), weak ($10^{-6}$--$10^{-4}$), or tiny/EH-dominated ($<10^{-6}$).

This diagnostic does not change the ranking, which remains by $\chiT$.  It changes the scientific interpretation.  An EH-only scalar model is a valid baseline.  A low-$\chi^2$ model with formally present but numerically irrelevant curvature corrections should not be advertised as an active modified-gravity discovery.  Conversely, a model with comparable $\chi^2$ and an active or moderate correction is a more interesting target for follow-up perturbation and structure-growth tests.

\section{Discussion}
\label{sec:discussion}

The current search suggests three lessons.  First, the combination of grammar validity and ghost-free FLRW filtering is essential.  Many covariant expressions are syntactically legal but physically unusable because they are higher derivative, trivial after coefficient relations, or impossible to convert into a regular IVP.  Second, the DES-SN plus DESI-DR2 BAO background likelihood is highly permissive toward scalar dark-energy histories close to an EH anchor.  This explains why simple $\phi_1$ sectors are repeatedly rediscovered across different proposal strategies.  Third, the Hamiltonian-type constraint is not an optional numerical detail.  It is the physical condition that selects initial data and prevents the Solver from accepting trajectories that solve only a rearranged RHS system while drifting away from the original diffeomorphism constraint.

The main limitation is that the validation is currently background-level.  Passing the present pipeline means that a candidate admits a ghost-free FLRW branch, a preserved modified Friedmann constraint, a numerically stable expansion history, and a competitive distance-redshift fit.  It does not guarantee absence of scalar or tensor perturbation ghosts, positive sound speeds, consistency with gravitational-wave propagation, viable screening, or agreement with growth-of-structure data.  These should be added as future deterministic validation layers.

The autonomous-search aspect is valuable because it changes how failures are used.  A rejected theory is not discarded silently; its failure mode becomes information.  Constraint-root failures point to poor initial-data compatibility.  Constraint drift indicates a problematic branch or numerical instability.  Failure to determine $q_{ix}$ points to an unusable scalar dynamical subsystem.  Symbolic timeouts identify operator combinations that are too complex for the current analytic reduction.  These are physical diagnostics first and machine-learning signals second.

\section{Conclusion}
\label{sec:conclusion}

We have described a physics-centered autonomous framework for discovering two-scalar modified-gravity theories from cosmological background data.  The framework generates dimension-balanced covariant Lagrangians, filters them through ghost-free FLRW equations, identifies the lapse equation as a preserved modified Friedmann constraint using residual time-diffeomorphism invariance, builds dimensionless IVP systems, and scores accepted trajectories against DES-SN plus DESI-DR2 BAO.

Recent scans find competitive background fits with $\chiT\simeq1638$.  The best families are often EH-anchored scalar sectors, while some accepted theories contain moderate or active curvature corrections.  This is exactly the kind of output a useful autonomous theory-discovery system should provide: not a single overclaimed answer, but a ranked set of physically validated candidates, their failure modes, and a clear distinction between baseline scalar dark energy and genuinely active modified-gravity corrections.

The next step is to extend the same deterministic validation philosophy beyond background cosmology.  Perturbation stability, gravitational-wave constraints, local tests, and structure growth should become additional hard filters.  With those layers in place, autonomous proposal mechanisms can explore a much larger theory space while the physics pipeline remains the final arbiter of viability.

\begin{acknowledgments}
We thank the developers of the AI Modified Gravity v5.2.2-agent pipeline and the scientific Python ecosystem used in the numerical implementation.
\end{acknowledgments}

\bibliographystyle{apsrev4-2}
\bibliography{main2}

\end{document}
